% BANKS-SOURCE: pdf=236 print=226 kind=prose
The four-index epsilon, Hodge-duality, instanton, and anomaly formulas in this
section are the $D=4$ Euclidean specialization, with orientation
$\epsilon_{0123}=+1$.

There are two possible Lorentz-invariant, gauge-invariant, dimension-4 terms in
the $D=4$ Euclidean Lagrangian of pure non-abelian gauge theory. These are

% BANKS-SOURCE: pdf=236 print=226 kind=equation id=10.7
\begin{equation}
  \mathcal L_E=\frac{1}{2g^2}\operatorname{tr}F_{\mu\nu}^2,
  \label{eq:10.7}
\end{equation}
\footnote[9]{The trace is taken in the fundamental representation with the
convention $\operatorname{tr}(\lambda_a\lambda_b)=\tfrac12\delta_{ab}$.}

and

% BANKS-SOURCE: pdf=236 print=226 kind=equation id=10.8
\begin{equation}
  \mathcal L_{\theta,E}=\frac{\theta}{32\pi^2}
  \operatorname{tr}F_{\mu\nu}(\ast F)^{\mu\nu},
  \label{eq:10.8}
\end{equation}

where

% BANKS-SOURCE: pdf=236 print=226 kind=equation id=10.9
\begin{equation}
  (\ast F)_{\mu\nu}=\frac12\epsilon_{\mu\nu\lambda\kappa}F^{\lambda\kappa}.
  \label{eq:10.9}
\end{equation}

% BANKS-SOURCE: pdf=236 print=226 kind=prose
The second of these is actually a total derivative. This is easy to see for the
abelian case. For the non-abelian case, we use the language of matrix-valued
differential forms. The gauge potential $A=\ii A_\mu\dd x^\mu$ is a
Lie-algebra-valued 1-form. The field strength $F$ is given by

% BANKS-SOURCE: pdf=236 print=226 kind=display
\[
  F=\dd A-A^2.
\]

Note that $F$ and $A$ are defined to be anti-Hermitian. Furthermore, because of
the anti-symmetry of the multiplication of forms, $A^2$ involves only the
commutator of the matrices $A_\mu$. Finally

% BANKS-SOURCE: pdf=236 print=226 kind=display
\[
  \dd F=-\dd A\,A+A\,\dd A=[A,F].
\]

% BANKS-SOURCE: pdf=237 print=227 kind=prose
Introduce the Chern--Simons form

% BANKS-SOURCE: pdf=237 print=227 kind=display
\[
  C=\operatorname{tr}\left(A\,\dd A-\frac23A^3\right).
\]

Then

% BANKS-SOURCE: pdf=237 print=227 kind=display
\[
\begin{aligned}
  \dd C
  &=\operatorname{tr}\left[(\dd A)^2-\frac23
    (\dd A\,A^2-A\,\dd A\,A+A^2\,\dd A)\right]\\
  &=\operatorname{tr}[\dd A\,F-2\,\dd A\,A^2],
\end{aligned}
\]

where we have used cyclicity of the trace and graded anti-commutativity of
differential forms. Noting that $\operatorname{tr}A^4=-\operatorname{tr}A^4=0$,
we conclude that

% BANKS-SOURCE: pdf=237 print=227 kind=equation id=10.10
\begin{equation}
  \dd C=\operatorname{tr}F^2
  \label{eq:10.10}
\end{equation}
\BanksImplicitHook{I-CH10-013}

% BANKS-SOURCE: pdf=237 print=227 kind=prose
is proportional to the $\theta$ term in the Lagrangian. The second term in the
action is therefore an integral over the boundary of Euclidean space-time of the
Chern--Simons form. The condition of finite action is that $A\to U^\dagger\dd U$,
a pure gauge, on the boundary. The function $U$ defines a map from the 3-sphere
at infinity in $\mathbb R^4$ to the gauge group. Consider the case of $SU(2)$,
where the group manifold is itself $S^3$. There are obviously different
topological classes of map of the sphere into itself, which are multiple
wrappings. They are characterized by the winding number $n$, which is an
arbitrary integer. Thus, there are different classes of finite-action gauge-field
configurations on $\mathbb R^4$, classified by the winding number of the pure
gauge transformation, which they approach at infinity. The winding number $n$
is called the topological charge of the configuration. For more general gauge
groups it can be shown that every non-trivial map of $S^3$ into the group can be
deformed into a map into some $SU(2)$ subgroup, so the topological classification
is the same.

In this $D=4$ Euclidean theory, note the obvious inequalities

% BANKS-SOURCE: pdf=237 print=227 kind=equation id=10.11
\begin{equation}
  \int\dd^4x\,\operatorname{tr}(F_{\mu\nu}\pm\ast F_{\mu\nu})^2\geq0.
  \label{eq:10.11}
\end{equation}

% BANKS-SOURCE: pdf=237 print=227 kind=prose
On expanding out the square and collecting terms we see that this bounds the
action of any configuration from below,

% BANKS-SOURCE: pdf=237 print=227 kind=display
\[
  S\geq\frac{8\pi^2}{g^2}|n|,
\]

where $n$ is the topological charge. Equality is achieved only for self-dual
and anti-self-dual configurations

% BANKS-SOURCE: pdf=237 print=227 kind=equation id=10.12
\begin{equation}
  F=\pm\ast F.
  \label{eq:10.12}
\end{equation}
\BanksImplicitHook{I-CH10-014}

Solutions of these first-order equations are automatically solutions of the
Euclidean Yang--Mills equations, because they minimize the action for fixed
topological charge. A basis of (anti-)self-dual tensors in four dimensions is
given by the 't Hooft symbols

% BANKS-SOURCE: pdf=237 print=227 kind=equation id=10.13
\begin{equation}
  \eta^a{}_{\mu\nu}=\delta_{\mu0}\delta_{\nu a}
  -\delta_{\nu0}\delta_{\mu a}+\epsilon_{\mu\nu a},
  \label{eq:10.13}
\end{equation}

% BANKS-SOURCE: pdf=237 print=227 kind=equation id=10.14
\begin{equation}
  \bar\eta^a{}_{\mu\nu}=\delta_{\mu0}\delta_{\nu a}
  -\delta_{\nu0}\delta_{\mu a}-\epsilon_{\mu\nu a}.
  \label{eq:10.14}
\end{equation}

% BANKS-SOURCE: pdf=238 print=228 kind=prose
These appear in the product rule for Euclidean Weyl matrices

% BANKS-SOURCE: pdf=238 print=228 kind=equation id=10.15
\begin{equation}
  \sigma^\mu\sigma^\nu=\delta^{\mu\nu}+\eta^{\mu\nu a}\sigma_a,
  \qquad
  \bar\sigma^\mu\bar\sigma^\nu=\delta^{\mu\nu}
  +\bar\eta^{\mu\nu a}\sigma_a.
  \label{eq:10.15}
\end{equation}

It is natural to search for a minimal-action solution that is maximally
symmetric, via the ansatz

% BANKS-SOURCE: pdf=238 print=228 kind=equation id=10.16
\begin{equation}
  A_\mu=g_1^{-1}(\hat x)\partial_\mu g_1(\hat x)f(x^2).
  \label{eq:10.16}
\end{equation}

$g_1$ is the winding number 1 mapping of the 3-sphere into $SU(2)$,

% BANKS-SOURCE: pdf=238 print=228 kind=equation id=10.17
\begin{equation}
  g_1=\hat x_\mu\sigma^\mu,
  \label{eq:10.17}
\end{equation}

where $\hat x$ is the unit vector in the $x$ direction. In Problem 10.5 the
reader will show that this gauge configuration is invariant under simultaneous
space-time and gauge rotations, and that its field strength is self-dual if

% BANKS-SOURCE: pdf=238 print=228 kind=display
\[
  f=\frac{x^2}{x^2+\rho^2}.
\]

$\rho$ is an arbitrary positive parameter, called the scale size of the
instanton.

In order to find a tunneling interpretation for this Euclidean solution we
consider Yang--Mills theory on the Euclidean manifold $\mathbb R\times S^3$.
Note that this manifold is conformally equivalent to $\mathbb R^4$, so we can
use the conformal invariance of the classical Yang--Mills equations to construct
the instanton on this manifold from the solution we have just found. Using the
fact that $F\wedge F=\dd C$, we find that the topological charge of the
instanton for this manifold is equal to the difference of the winding numbers of
the pure gauge configurations, which it approaches as $t\to\pm\infty$.

% BANKS-SOURCE: pdf=238 print=228 kind=subsection id=10.6.1
\subsubsection{Hamiltonian formulation of Yang--Mills theory in temporal gauge}

% BANKS-SOURCE: pdf=238 print=228 kind=prose
It is always possible to use the gauge freedom of a non-abelian gauge theory to
set $A^a{}_{0}=0$. When this is done, the non-abelian electric field is just

% BANKS-SOURCE: pdf=238 print=228 kind=display
\[
  \mathbf E^{\,a}=\frac1{g^2}\partial_0\mathbf A^{\,a}.
\]

The Lagrangian has a standard canonical form and can be quantized in a
straightforward manner, with $E_i{}^a$ and $A_i{}^a$ as canonical conjugates.
However, we must also impose the condition that follows from varying the
original action with respect to $A^a{}_{0}$. It is easy to see that this is

% BANKS-SOURCE: pdf=238 print=228 kind=display
\[
  G^a(\mathbf x)=D_i{}^{ab}E_i{}^b-\rho^a=0.
\]

This is the non-abelian form of Gauss' law, and $\rho^a$ is the time component
of the Noether current for matter fields, which follows from the global $G$
symmetry. This

% BANKS-SOURCE: pdf=239 print=229 kind=prose
condition should be understood as follows. The system with $A^a{}_{0}=0$ is
invariant under *time-independent* gauge transformations, and

% BANKS-SOURCE: pdf=239 print=229 kind=display
% CONVENTION-SCOPE: fixed D=4 reason=temporal-gauge spatial generator in D=4 Lorentzian Yang--Mills
\[
  G(\omega)=\int\dd^3\mathbf x\,\omega^a(\mathbf x)G^a(\mathbf x)
\]

is the generator of these transformations when $\omega^a$ vanishes sufficiently
rapidly at infinity that integrations by parts are permitted.

It is easy to verify that each $G(\omega)$ commutes with the Hamiltonian, so that
we have an infinite-dimensional symmetry group of the mechanical system we have
defined. The statement that $G(\omega)=0$ is imposed on the allowed physical
states of this theory. It is somewhat analogous to the restriction to
BRST-invariant states in covariant quantization.

An interesting twist on this formalism occurs if there is a topological term in
the Yang--Mills action, and $G$ is broken down to $H=U(1)$ (or a product of
$U(1)$ groups). The topological term does not affect the equations of motion, but
it does change the canonical momentum, and thus the generator of gauge
transformations. In electrodynamics, it is easy to see that the constraint
equation is now

% BANKS-SOURCE: pdf=239 print=229 kind=display
\[
  \partial_iE_i-\frac{\theta}{2\pi}\partial_iB_i-\rho=0,
\]

where $\rho$ is the electric charge density of matter fields. We see that, in the
presence of $\theta$, magnetic monopoles will carry (generally irrational)
electric charges \cite{banks-ref-140}.
\BanksImplicitHook{I-CH10-015}

In the classical approximation, ground states of the theory are found by looking
for minima of the energy. The classical states of minimum energy are static,
pure gauge-field configurations. When space has the topology $S^3$ (equivalently,
when it has the topology $\mathbb R^3$ but we impose falloff conditions at
infinity), the space of such field configurations is classified by the
topological winding number $n$. The instanton represents a tunneling process
between two classical vacuum states with winding numbers differing by one. This
is closely analogous to the tunneling between different wells of an infinite
periodic potential in one dimension. The true semi-classical vacuum state in
that case is a Bloch wave, characterized by a Bloch momentum, which runs between
$0$ and $2\pi$. In an analogous fashion, semi-classical Yang--Mills theory would
appear to have a continuous set of vacua characterized by a parameter $\theta$
that is periodic with period $2\pi$. In fact this is true, even though the
semi-classical approximation is valid only for Green functions at short
distances, in the confining phase of the theory. The parameter $\theta$ is just
the coefficient of the $F\wedge F$ term in the Lagrangian, which we introduced
above. A simple example of the correspondence between vacuum parameters and
terms in the Lagrangian is treated in Problem 10.2.

% BANKS-SOURCE: pdf=239 print=229 kind=subsection id=10.6.2
\subsubsection{Bosonic zero modes}

% BANKS-SOURCE: pdf=239 print=229 kind=prose
There is also an anti-instanton solution of the anti-self-duality equations in
which the winding-number-1 gauge transformation is replaced by the
transformation of opposite winding number. Atiyah and Hitchin were able to find
all solutions of the self-duality

% BANKS-SOURCE: pdf=240 print=230 kind=prose
equations \cite{banks-ref-141}. These exact solutions are useful primarily in
supersymmetric gauge theories, where they allow us to calculate certain
correlation functions exactly. Their study goes beyond the scope of this book.

For our purposes, it is more interesting to look at approximate solutions of the
second-order equations consisting of a dilute gas of widely separated instantons
and anti-instantons. In this context, widely separated means separations large
compared with all of the parameters $\rho_i$ of the instantons and
anti-instantons. The existence of the parameter $\rho$ is a consequence of the
classical conformal invariance of the Yang--Mills equations. Given any solution
with finite action, we can scale the coordinates to get a new solution with the
same action. Note that the same is not true for finite special conformal
transformations, because they map finite points to infinity. When we consider
small fluctuations of the instanton, these symmetries show up as zero modes of
the fluctuation operator. The zero modes corresponding to special conformal
transformations are not normalizable and the corresponding directions in field
space are not included in the functional integral.

The procedure for dealing with the scale zero mode is similar to that for
translations. One extracts a scale collective coordinate and does the integral
over it exactly. The classical measure for this integration is the scale-invariant
measure

% BANKS-SOURCE: pdf=240 print=230 kind=display
\[
  \frac{\dd\rho}{\rho},
\]

but there are generally quantum corrections to this. In Yang--Mills theory, the
dominant effect is the replacement of $g^2$ in the instanton action by the running
value of the renormalized coupling $g^2(\rho)$ \cite{banks-ref-142}. In
asymptotically free theories $g$ goes to zero for zero scale size, which tends to
make the integral over $\rho$ converge near 0. On the other hand, this leads to
an enhanced divergence at large $\rho$. Instanton calculations of the
short-distance behavior of Green functions are under control, because the $\rho$
integrals get cut off at scales of order the distance between operators. However,
instanton technology, just like perturbation theory, fails to capture the
large-distance behavior of asymptotically free theories.

In a purely bosonic gauge theory, the result of the dilute-gas computation of
Green functions is

% BANKS-SOURCE: pdf=240 print=230 kind=equation id=10.18
\begin{equation}
\begin{aligned}
  \langle O_1(\mathbf x_1)\ldots O_n(\mathbf x_n)\rangle_c
  ={}&D^{-1/2}\int\dd^4\mathbf a\int\frac{\dd\rho}{\rho}\int\dd\Omega\\
  &\times g^{-(5+d_G-3)}(\rho)\ee^{-\frac{8\pi^2}{g^2(\rho)}}
  O_1^{(c)}(\mathbf x_1+\mathbf a)\ldots
  O_n^{(c)}(\mathbf x_n+\mathbf a).
  \label{eq:10.18}
\end{aligned}
\end{equation}

% BANKS-SOURCE: pdf=240 print=230 kind=prose
The integrals are over the position, scale size, and embedding of the instanton
configuration in the gauge group. $O_k^{(c)}$ is the classical value of the
operator $O_k$ in the instanton configuration of fixed position, size, and gauge
orientation. The factor $D^{-1/2}$ is the determinant computed from the
non-zero modes of fluctuation around the instanton. It has been computed by 't
Hooft \cite{banks-ref-142}. The inverse powers of $g$ come from the normalization
of the zero mode integrals, as explained above.

% BANKS-SOURCE: pdf=241 print=231 kind=subsection id=10.6.3
\subsubsection{Fermion zero modes and anomalies}

% BANKS-SOURCE: pdf=241 print=231 kind=prose
In gauge theories with fermions, an extremely interesting phenomenon occurs in
connection with zero modes of the massless Dirac operator in an instanton
background. The chirality statements use $D=4$, where the matrix $\gamma_5$
exists and anticommutes with the Dirac matrices. Consider the eigenvalue
equation for the Euclidean Dirac operator:

% BANKS-SOURCE: pdf=241 print=231 kind=equation id=10.19
\begin{equation}
  \ii\slashed D\psi=\lambda\psi.
  \label{eq:10.19}
\end{equation}

$\gamma_5\psi$ is then an eigenfunction as well, with eigenvalue $-\lambda$.
Now consider modes with $\lambda=0$, and the way in which they behave under
continuous deformations of the Yang--Mills background. In general, we would
expect such a deformation to change the eigenvalue, but we can only change pairs
of eigenvalues away from zero.

For $\lambda=0$, the eigenfunctions may be chosen to have a definite chirality,
$\gamma_5\psi=\pm\psi$, because $\gamma_5$ anti-commutes with the Dirac
operator and preserves the zero eigenspace. This is not true for $\lambda\ne0$,
since we have seen that multiplication by $\gamma_5$ reverses the sign of the
eigenvalue. As a consequence, the wave function of a non-zero eigenvalue is a
linear combination of the two different chiralities. Thus, left- and
right-handed zero modes must be lifted in pairs.

This means that the difference between the number of left- and right-handed zero
modes (the so-called *index* of the Dirac operator) is a topological invariant of
the gauge field, unchanged by continuous deformations. We might be tempted to
think that this is related to the topological invariant we have just been
discussing. A famous mathematical theorem, due to Atiyah and Singer
\cite{banks-ref-143,banks-ref-144,banks-ref-145,banks-ref-146}, tells us that
this is the case. A physics proof of this equation follows from the anomaly
equation for the axial $U(1)$ current of the massless fermions

% BANKS-SOURCE: pdf=241 print=231 kind=equation id=10.20
\begin{equation}
  \partial_\mu\bar\psi\gamma^\mu\gamma_5\psi
  =\frac{1}{32\pi^2}F^a{}_{\mu\nu}(\ast F)^{\mu\nu}{}_{a}.
  \label{eq:10.20}
\end{equation}

Consider the Euclidean fermion functional integral in a background gauge field
with finite action and topological charge. The functional integral of the
left-hand side of the anomaly equation is

% BANKS-SOURCE: pdf=241 print=231 kind=equation id=10.21
\begin{equation}
  \sum_n\partial_\mu\bar\psi_n\gamma^\mu\gamma_5\psi_n,
  \label{eq:10.21}
\end{equation}

where the sum is over all normalized eigenfunctions of the Dirac operator,

% BANKS-SOURCE: pdf=241 print=231 kind=equation id=10.22
\begin{equation}
  \slashed D\psi_n=\lambda_n\psi_n.
  \label{eq:10.22}
\end{equation}

Since the gauge potential falls off at infinity, the finite-$\lambda_n$
eigenfunctions fall off exponentially at infinity. Thus, if we integrate (10.21)
over all of Euclidean space, only the zero modes contribute, and the integral
just counts the difference between the numbers of right- and left-handed zero
modes. On the other hand, the right-hand side integrates to the topological
charge. So we find that the topological charge is in fact equal to the difference
between the numbers of right- and left-handed zero modes.

Actually, this calculation is valid only when the fermions are in the fundamental
representation. For a general fermion representation, $R$, if we choose a
representative of the topological class for which the gauge field sits in an
$SU(2)$ subgroup, then $R$

% BANKS-SOURCE: pdf=242 print=232 kind=prose
will break up into a direct sum of spin-$j$ representations of $SU(2)$. The
anomaly of the spin-$j$ representation is
$2\sum_{m=-j}^{j}m^2$ times larger than that of the spin-$\tfrac12$
representation, so we can use this formula to count the number of zero modes for
any representation.

Bosonic zero modes give apparently infinite contributions to the path integral.
We have to integrate over the corresponding collective coordinates exactly in
order to eliminate or interpret the infinity. By contrast, fermion zero modes
tell us that the single-instanton contribution to the partition function
vanishes. Instead, the single instanton contributes to the expectation value of
operators that can ``absorb'' the zero modes. The simplest operator, with
minimal number of fields, which can do the job, is called the 't Hooft operator
associated with the instanton. The dilute-instanton-gas approximation for
theories with fermion zero modes corresponds to the lowest-order perturbation
theory of a modified theory in which the 't Hooft operator is added to the
Lagrangian with coefficient $D\ee^{-S_I}$, where $D$ includes the determinant of
non-zero modes, and inverse coupling factors for bosonic zero modes. These
expressions must be integrated over bosonic collective coordinates.

The anomaly equation tells us that the $U(1)$ axial symmetry is broken by
instantons. This shows up in the structure of the 't Hooft operator, which
explicitly violates the symmetry. In general, if there are $k$ simple factors of
the gauge group, we may expect to have $k$ independent symmetries of the
classical Lagrangian, which are broken by these non-perturbative effects.

% Keep the printed Chapter 10 footnote sequence aligned for the soliton section.
\setcounter{footnote}{9}
